% !TEX root = ../m1-20-21.tex
\chapter{Șiruri și serii de funcții}

\section{Șiruri de funcții}

\subsection{Convergență punctuală și convergență uniformă}

Fie $ (f_n)_n : \RR \to \RR $ un șir (o familie) de funcții, adică
pentru fiecare $ n \in \NN $, avem cîte o funcție $ f_n : \RR \to \RR $.
Putem gîndi aceste șiruri de funcții ca pe niște generalizări ale
șirurilor de numere reale. Dacă în cazul acela, regula generală era fixată,
dată de $ x_n = x(n) $, în cazul șirurilor de funcții, pentru fiecare indice
$ n $, putem avea cîte o regulă diferită, dată de o funcție $ f_n $.

Dacă șirurile de numere reale au drept limită un număr real și astfel
putem spune că un șir de numere reale aproximează un număr real, în cazul
șirurilor de funcții, limita va fi, în general, o funcție. Deci putem spune
că aceste șiruri aproximează funcții. Însă noțiunea de convergență în cazul
șirurilor de funcții este ceva mai subtilă.

\begin{definition}\label{def:conv-simpl}\index{convergență!simplă}
  Fie $ (f_n)_n : D \seq \RR \to \RR $ un șir de funcții.

  Spunem că șirul $ (f_n)_n $ \emph{converge punctual (simplu)} la
  funcția $ f $ dacă are loc:
  \[
    \lim_{n \to \infty} f_n(x) = f(x), \quad \forall x \in D.
  \]
  Pe scurt, putem nota aceasta prin $ f_n \xrar{PC} f $ sau $ f_n \xrar{s} f $,
  iar funcția $ f $ se va numi \emph{limita punctuală} a șirului $ (f_n) $.
\end{definition}

Acest tip de convergență este foarte puternic, deoarece este definit destul
de grosier, astfel că avem nevoie de rafinări ale conceptului.

Celălalt tip de convergență este definit mai jos.

\begin{definition}\label{def:conv-unif}\index{convergență!uniformă}
  În condițiile și cu notațiile de mai sus, spunem că șirul $ (f_n) $ este
  \emph{uniform convergent} la funcția $ f $ dacă:
  \[
    \forall \varepsilon > 0, \exists N_\varepsilon > 0 \text{ a. î. } %
    |f_n(x) - f(x)| < \varepsilon, \forall n \geq N_\varepsilon, \forall x \in D.
  \]
  Notația va fi $ f_n \xrar{UC} f $ sau $ f_n \xrar{u} f $.
\end{definition}

Cu alte cuvinte, șirul $ (f_n) $ se păstrează arbitrar de aproape de funcția limită,
de la un anumit rang încolo.

Ca și în cazul șirurilor de numere, se va folosi mai des în exerciții o teoremă
de caracterizare.

\begin{proposition}\label{prop:conv-u}
  În condițiile și cu notațiile de mai sus, are loc:
  \[
    f_n \xrar{u} f \Leftrightarrow \lim_{n \to \infty} \sup_{x \in D} |f_n(x) - f(x)| = 0.
  \]
\end{proposition}

Se poate vedea că proprietatea de convergență uniformă o implică pe cea de 
convergență punctuală, însă reciproca este falsă.

De exemplu, fie $ f_n : [0, 1] \to \RR, f_n(x) = x^n, n \geq 1 $.

Atunci se poate vedea imediat că:
\[
  \lim_{n \to \infty} f_n(x) = %
  \begin{cases}
    0, & 0 \leq x < 1 \\
    1, & x = 1
  \end{cases}.
\]
De aici rezultă că $ f_n \xrar{s} f $, unde funcția $ f $ este definită pe cazuri
mai sus.

Dar un calcul simplu arată că:
\[
  \lim_{n \to \infty} \sup_{x \in [0, 1]} |f_n(x) - f(x)| = 1 \neq 0,
\]
deci șirul nu este uniform convergent.

%%%%%%%%%%%%%%%%%%%%%%%%%%%%%%%%%%%%%%%%%%%%%%%%%%%%%%%%%%%%%%%%%%%%%% 

\subsection{Transferul proprietăților}

În unele situații, putem verifica proprietățile de convergență simplă și
mai ales uniformă prin transferul proprietăților. Cu alte cuvinte, vom ști
ce proprietăți analitice (continuitate, integrabilitate) ale funcțiilor se
păstrează prin convergență și aceasta ne va da condiții necesare pentru
a studia acel tip de convergență.

\begin{theorem}[Transfer de continuitate]\label{thm:transfer-cont}\index{transfer!de continuitate}
  Fie $ f_n : D \seq \RR \to \RR $ un șir de funcții.

  Dacă fiecare $ f_n $ este funcție continuă, iar șirul $ f_n $ converge uniform
  la funcția $ f $, atunci $ f $ este continuă.
\end{theorem}

Această teoremă ne ajută să verificăm convergență uniformă astfel: dacă știm
(am demonstrat) că fiecare termen $ f_n $ este funcție continuă, iar
funcția $ f $, obținută prin convergența \emph{punctuală} a șirului (singurul
candidat posibil și pentru convergența uniformă!) nu este o funcție continuă,
atunci știm sigur că șirul nu converge uniform la $ f $.

\begin{theorem}[Integrare termen cu termen]\label{thm:transf-int}\index{transfer!de integrabilitate}
  Fie $ f_n : D \seq \RR \to \RR $ un șir de funcții și $ f : D \seq \RR \to \RR $ o
  funcție.

  Dacă $ f_n \xrar{u} f $, atunci are loc proprietatea de integrare termen
  cu termen, adică:
  \[
    \lim_{n \to \infty} \int_a^b f_n(x) dx = \int_a^b f(x) dx.
  \]
\end{theorem}

Cu alte cuvinte, dacă știm că are loc convergența uniformă, atunci limita și integrala
pot comuta.

\begin{theorem}[Derivare termen cu termen]\label{thm:transf-der}\index{transfer!de derivabilitate}
  Fie $ f_n : D \seq \RR \to \RR $ un șir de funcții și $ f : D \seq \RR \to \RR $ o funcție.

  Presupunem că funcțiile $ f_n $ sînt derivabile, pentru orice $ n \in \NN $. Dacă
  $ f_n \xrar{s} f $ și dacă există o funcție $ g : D \seq \RR \to \RR $ astfel încît
  $ f_n' \xrar{u} g $, atunci $ f $ este derivabilă și $ f' = g $.
\end{theorem}

Această proprietate este similară celei de integrabilitate și ne arată că, dacă
are loc convergența punctuală a șirului funcțiilor și uniformă a șirului derivatelor,
atunci limita și derivata pot comuta.


%%%%%%%%%%%%%%%%%%%%%%%%%%%%%%%%%%%%%%%%%%%%%%%%%%%%%%%%%%%%%%%%%%%%%% 

\section{Serii de funcții}
\index{serii!de funcții}

Pasul următor este să studiem serii de funcții. Pentru convergența acestora,
avem un singur criteriu de utilizat.

\begin{theorem}[Weierstrass]\label{thm:weierstrass}\index{criteriul!Weierstrass}
  Fie $ (a_n) $ un șir cu termeni pozitivi, $ \sum f_n(x) $ o serie de funcții, cu
  fiecare $ f_n : D \seq \RR \to \RR $, astfel încît $ |f_n(x)| \leq a_n $
  pentru orice $ x \in D $, iar seria $ \sum a_n $ este convergentă,
  atunci seria $ \sum f_n $ converge uniform.
\end{theorem}

Transferul proprietăților se formulează pe scurt astfel (vom păstra notațiile
și contextul din teorema lui Weierstrass de mai sus):
\begin{itemize}
\item \emph{Transfer de continuitate:} Dacă $ f_n $ sînt funcții continue,
  iar seria $ \sum f_n $ converge uniform la $ f $, atunci funcția $ f $
  este continuă.
\item \emph{Integrare termen cu termen:} Dacă seria $ \sum f_n $ converge
  uniform la $ f $, atunci $ f $ este integrabilă și are loc:
  \[
    \int_a^b \sum_n f_n(x) dx = \sum_n \int_a^b f_n(x) dx,
  \]
  pentru orice interval $ [a, b] \seq D \seq \RR $ din domeniul de
  definiție al funcțiilor $ f_n $.
\item \emph{Derivare termen cu termen}: Presupunem că toate funcțiile
  $ f_n $ sînt derivabile. Dacă seria $ \sum f_n $ converge punctual
  la $ f $ și dacă există $ g : D \to \RR $ astfel încît $ \sum f'_n $
  converge uniform la $ g $, atunci $ f $ este derivabilă și $ f' = g $.
\end{itemize}

Toate aceste proprietăți, însă, vor fi utile într-un caz particular de serii de
funcții, anume acela al seriilor de puteri.


%%%%%%%%%%%%%%%%%%%%%%%%%%%%%%%%%%%%%%%%%%%%%%%%%%%%%%%%%%%%%%%%%%%%%% 

\section{Serii de puteri}
\index{serii!de puteri}

Pornim cu o serie de funcții $ \sum f_n $, dar în care fiecare termen $ f_n $
este o funcție de tip polinomial. Astfel că, de fapt, seriile de puteri
pot fi scrise în general sub forma $ \sum a_n(x - a)^n $, cu $ a_n, a \in \RR $,
caz în care se numesc \emph{serii de puteri centrate în $ a $}, definite de
șirul $ (a_n) $.

În majoritatea cazurilor, vom lucra cu serii de puteri centrate în origine,
care se vor scrie, în general, $ \sum a_n x^n $.

Toate rezultatele de la serii de funcții sînt valabile și cu demonstrații imediate,
dat fiind că funcțiile polinomiale sînt ușor de studiat. Rezultatele specifice
sînt următoarele.

\begin{theorem}[Abel]\label{thm:abel}\index{teorema!Abel}
  Fie $ \sum a_n (x - a)^n $ o serie de puteri centrată în $ a \in \RR $. Atunci
  există un număr $ 0 \leq R \leq \infty $ astfel încît:
  \begin{itemize}
  \item Seria este absolut convergentă pe intervalul $ (a - R, a + R) $;
  \item Seria este divergentă pentru $ |x| > R $;
  \item Seria este uniform convergentă pentru $ [-r, r] $, pentru $ 0 < r < R $.
  \end{itemize}
  Numărul $ R $ se numește \emph{raza de convergență a seriei}, iar intervalul
  $ (a - R, a + R) $ se numește \emph{intervalul de convergență}.
\end{theorem}

Problema centrală acum va fi calculul razei de convergență, care se poate
face cu unul dintre criteriile de mai jos.

\begin{theorem}[Cauchy-Hadamard]\label{thm:cauchy-had}\index{teorema!Cauchy-Hadamard}
  Fie $ \sum a_n (x - a)^n $ o serie de puteri. Fie $ R $ raza sa de convergență
  și fie $ \omega = \lim\sup \sqrt[n]{|a_n|} $. Atunci:
  \begin{itemize}
  \item $ R = \omega^{-1} $, pentru $ 0 < \omega < \infty $;
  \item $ R = 0 $, dacă $ \omega = \infty $;
  \item $ R = \infty $, dacă $ \omega = 0 $.
  \end{itemize}
\end{theorem}

\begin{theorem}\label{thm:raza-conv}
  În condițiile și cu notațiile de mai sus, raza de convergență se poate calcula și 
  cu formula:
  \[
    R = \lim_{n \to \infty} \left| \dfrac{a_n}{a_{n+1}} \right|.
  \]
\end{theorem}

\begin{remark}\label{rk:transfer-serii-puteri}
  Remarcăm că proprietățile de transfer sînt imediate în cazul seriilor de puteri.
  Rezultă că, dacă $ \sum a_n (x - a)^n $ este o serie de puteri convergentă la
  $ S(x) $, atunci:
  \begin{itemize}
  \item Seria derivatelor, $ \sum n a_n(x - a)^{n-1} $, are aceeași rază de
    convergență cu seria inițială și are suma $ S'(x) $;
  \item Seria primitivelor, $ \sum \dfrac{a_n}{n+1} (x - a)^{n+1} $ are
    aceeași rază de convergență cu seria inițială, iar suma sa este
    o primitivă a funcției $ S $.
  \end{itemize}
\end{remark}

În exerciții, se va folosi foarte des seria geometrică, împreună cu derivata
și o primitivă a ei. Astfel, pentru cazul convergent ($ | q | < 1 $), avem:
\begin{align*}
  \sum aq^n &= \dfrac{a}{1 - q} \Rightarrow\\
  \sum a\dfrac{q^{n + 1}}{n + 1} &= -a\ln(1 - q) \\
  \sum anq^{n-1} &= \dfrac{-a}{(1 - q)^2}.
\end{align*}



%%%%%%%%%%%%%%%%%%%%%%%%%%%%%%%%%%%%%%%%%%%%%%%%%%%%%%%%%%%%%%%%%%%%%% 

\section{Polinomul Taylor și seria Taylor}

Orice funcție cu proprietăți analitice \qq{bune} poate fi aproximată cu un polinom.

\begin{definition}\label{def:pol-taylor}\index{polinom!Taylor}
  Fie $ D \seq \RR $ un interval deschis și $ f : D \to \RR $ o funcție de clasă
  $ \kal{C}^m(D) $, adică este continuă și are derivata continuă, pînă la a $ m $-a
  iterație.

  Pentru orice $ a \in D $ definim \emph{polinomul Taylor} de grad $ n \leq m $
  asociat funcției $ f $ în punctul $ a $ prin:
  \[
    T_{n, f, a} = \sum_{k = 0}^n \dfrac{f^{(k)}(a)}{k!} (x - a)^k.
  \]
  Restul (eroarea de aproximare), numit și \emph{restul Lagrange} este definit
  prin:
  \[
    R_{n, f, a} = f(x) - T_{n, f, a}(x).
  \]
\end{definition}

Pentru cazul particular al polinomului Taylor de grad 1, acesta se mai numește
\emph{aproximația liniară} a funcției $ f $, iar pentru gradul 2,
\emph{aproximația pătratică}.

Evident, acest polinom poate fi folosit mai departe pentru a studia \emph{seria} Taylor
asociată unei funcții. Formal, conceptul este definit mai jos.

\begin{theorem}[Seria Taylor]\label{thm:seria-taylor}\index{serii!Taylor}
  Fie $ a < b $ și $ f \in \kal{C}^\infty([a, b]) $, astfel încît să existe
  $ M > 0 $ cu $ |f^{(n)}(x)| \leq M $, pentru orice $ n \in \NN, x \in [a, b] $.

  Atunci pentru orice $ x_0 \in (a, b) $, se definește seria Taylor a lui $ f $ în
  jurul punctului $ x_0 $, care va fi uniform convergentă pe $ [a, b] $, iar
  suma sa este funcția $ f $, adică avem:
  \[
    f(x) = \sum_{n \geq 0} \dfrac{f^{(n)}(x_0)}{n!} (x - x_0)^n, \quad \forall x \in [a, b].
  \]
  Pentru cazul particular $ x_0 = 0 $, seria se numește \emph{Maclaurin.}\index{serii!Maclaurin}
\end{theorem}

\begin{remark}\label{rk:mac}
  Dacă nu se specifică punctul în jurul căruia să se facă dezvoltarea în serie Taylor
  pentru funcția $ f $, vom presupune că lucrăm cu serii Maclaurin.
\end{remark}

\begin{remark}\label{rk:aprox}
  Conform proprietăților de aproximare a sumelor seriilor convergente din seminarul
  anterior, știm că eroarea de aproximare (restul Lagrange, în acest caz) este
  comparabilă cu primul termen neglijat din seria Taylor (echivalent, cu următorul
  termen din polinomul Taylor).
\end{remark}

\textbf{Seriile Taylor uzuale} pentru funcțiile elementare, împreună cu domeniile
de convergență, sînt date mai jos.
\begin{align*}
  e^x &= \sum_{n \geq 0} \dfrac{1}{n!} x^n, x \in \RR \\
  \dfrac{1}{1 - x} &= \sum_{n \geq 0} x^n, |x| < 1 \\
  \dfrac{1}{1 + x} &= \sum_{n \geq 0} (-1)^n x^n, |x| < 1 \\
  \cos x &= \sum_{n \geq 0} \dfrac{(-1)^n}{(2n)!} x^{2n}, x \in \RR \\
  \sin x &= \sum_{n \geq 0} \dfrac{(-1)^n}{(2n + 1)!} x^{2n + 1}, x \in \RR \\
  (1 + x)^\alpha &= \sum_{n \geq 0} \dfrac{\alpha(\alpha - 1)(\alpha - 2) \cdots (\alpha - n + 1)}{n!} x^n,%
                   |x| < 1, \alpha \in \RR \\
  \arctan x &= \sum_{n \geq 0} \dfrac{(-1)^n}{2n + 1} x^{2n + 1}, | x | \leq 1.
\end{align*}

Seriile pentru alte funcții se pot obține fie prin calcul direct, fie
prin derivare sau integrare termen cu termen a seriilor de mai sus.

%%%%%%%%%%%%%%%%%%%%%%%%%%%%%%%%%%%%%%%%%%%%%%%%%%%%%%%%%%%%%%%%%%%%%% 

\section{Exerciții}

1. Să se studieze convergența punctuală și uniformă a șirurilor de funcții:
\begin{enumerate}[(a)]
\item $ f_n : (0, 1) \to \RR, f_n(x) = \dfrac{1}{nx + 1}, n \geq 0 $;
\item $ f_n : [0, 1] \to \RR, f_n(x) = x^n - x^{2n}, n \geq 0 $;
\item $ f_n : \RR \to \RR, f_n(x) = \sqrt{x^2 + \dfrac{1}{n^2}}, n > 0 $;
\item $ f_n : [-1,1] \to \RR, f_n(x) = \dfrac{x}{1 + nx^2} $;
\item $ f_n : (-1, 1) \to \RR, f_n(x) = \dfrac{1 - x^n}{1 - x} $;
\item $ f_n : [0, 1] \to \RR, f_n(x) = \dfrac{2nx}{1 + n^2 x^2} $;
\item $ f_n : \RR_+ \to \RR, f_n(x) = \dfrac{x + n}{x + n + 1} $;
\item $ f_n : \RR_+ \to \RR, f_n(x) = \dfrac{x}{1 + nx^2} $.
\end{enumerate}

\vspace{1cm}

2. Să se arate că șirul de funcții dat de:
\[
  f_n : \RR \to \RR, \quad f_n(x) = \frac{1}{n} \arctan x^n
\]
converge uniform pe $ \RR $, dar:
\[
  \Big( \lim_{n \to \infty} f_n(x) \Big)'_{x = 1} \neq \lim_{n \to \infty} f'_n(1).
\]

Rezultatele diferă deoarece șirul derivatelor nu converge uniform pe $ \RR $.

\vspace{1cm}

3. Să se arate că șirul de funcții dat de:
\[
  f_n : [0, 1] \to \RR, \quad f_n(x) = nxe^{-nx^2}
\]
este convergent, dar:
\[
  \lim_{n \to \infty} \int_0^1 f_n(x) dx \neq \int_0^1 \lim_{n \to \infty} f_n(x) dx.
\]
Rezultatul se explică prin faptul că șirul nu este uniform convergent.
De exemplu, pentru $ x_n = \dfrac{1}{n} $, avem $ f_n(x_n) \to 1 $, dar,
în general, $ f_n(x) \to 0 $.

\vspace{1cm}

4. Să se dezvolte următoarele funcții în serie Maclaurin, precizînd și
domeniul de convergență:
\begin{enumerate}[(a)]
\item $ f(x) = e^x $;
\item $ f(x) = \sin x $;
\item $ f(x) = \cos x $;
\item $ f(x) = (1 + x)^\alpha, \alpha \in \RR $;
\item $ f(x) = \dfrac{1}{1 + x} $;
\item $ f(x) = \ln(1 + x) $;
\item $ f(x) = \arctan x $;
\item $ f(x) = \ln(1 + 5x)$;
\item $ f(x) = 3 \ln(2 + 3x) $.
\end{enumerate}

%%%%%%%%%%%%%%%%%%%%%%%%%%%%%%%%%%%%%%%%%%%%%%%%%%%%%%%%%%%%%%%%%%%%%% 

\vspace{1cm}

5. Să se calculeze raza de convergență și mulțimea de convergență pentru
următoarele serii de puteri:
\begin{enumerate}[(a)]
\item $ \sum_{n \geq 0} x^n $;
\item $ \sum_{n \geq 1} n^n x^n $;
\item $ \sum_{n \geq 1} (-1)^{n+1} \dfrac{x^n}{n} $;
\item $ \sum_{n \geq 1} \dfrac{n^nx^n}{n!} $;
\item $ \sum \dfrac{(x - 1)^{2n}}{n \cdot 9^n} $;
\item $ \sum \dfrac{(x + 3)^n}{n^2} $;
\item $ \sum \dfrac{n + 2}{n^2 + 1} (x - 2)^n $;
\item $ \sum \dfrac{n + 1}{\sqrt{n^4 + n^3 + 1}} \left( \dfrac{x + 1}{2x + 3} \right)^n $.
\end{enumerate}

\vspace{1cm}

6. Găsiți mulțimea de convergență și suma seriilor:
\begin{enumerate}[(a)]
\item $ \sum_{n \geq 0} (-1)^n \dfrac{x^{2n+1}}{2n + 1} $;
\item $ \sum \dfrac{2^n}{2n + 1} x^n $;
\item $ \sum \dfrac{n}{n + 1} x^n $.
\end{enumerate}

\vspace{1cm}

7. Să se calculeze cu o eroare mai mică decît $ 10^{-3}$ integralele:
\begin{enumerate}[(a)]
\item $ \ds\int_0^{\frac{1}{2}} \dfrac{\sin x}{x} dx $;
\item $ \ds\int_0^{\frac{1}{2}} \dfrac{\ln(1 + x)}{x} dx $;
\item $ \ds\int_0^1 e^{-x^2} dx $.
\end{enumerate}

\vspace{1cm}

8. Să se calculeze polinomul Taylor de grad 3 în jurul originii
pentru funcțiile:
\begin{enumerate}[(a)]
\item $ f(x) = 3 \ln(2 + x) $;
\item $ f(x) = \arctan x $;
\item $ f(x) = \sqrt{1 + 2x} $.
\end{enumerate}

\vspace{1cm}

9. Găsiți aproximarea liniară și pătratică a funcțiilor:
\begin{enumerate}[(a)]
\item $ f(x) = \sqrt[3]{x} $;
\item $ f(x) = \sin(\cos x) $;
\item $ f(x) = e^{\sin x} $;
\item $ f(x) = \arcsin x $.
\end{enumerate}

\vspace{1cm}

10. Folosind seria Taylor, aproximați cu o eroare mai mică decît $ 10^{-3} $
numerele:
\begin{enumerate}[(a)]
\item $ \sqrt[3]{65} $;
\item $ \sin 32 $;
\item $ \arctan \dfrac{1}{2} $;
\item $ e^{-0,2} $;
\item $ \ln 1,1 $;
\item $ \ln 4 $;
\item $ \ln 5 $.
\end{enumerate}

\emph{Indicație:} Atenție la domeniile de convergență!

\vspace{1cm}

11*. Arătați că seriile numerice de mai jos sînt convergente și calculați
sumele lor, folosind serii de puteri:
\begin{enumerate}[(a)]
\item $ \ds\sum_{n \geq 0} \dfrac{(-1)^n}{3n + 1} $;
\item $ \ds\sum_{n \geq 0} \dfrac{(n+1)^2}{n!} $;
\item $ \ds\sum_{n \geq 1} \dfrac{n^2(3^n - 2^n)}{6^n} $.
\end{enumerate}

\emph{Indicații:}
\begin{enumerate}[(a)]
\item Seria satisface criteriul lui Leibniz, deci este convergentă.

  Pentru a găsi suma, pornim cu seria de puteri $ \sum (-1)^n \dfrac{x^{3n+1}}{3n+1} $.

  Intervalul de convergență este $ (-1, 1) $, iar pentru $ x = 1 $, avem
  seria dată.

  Fie $ f $ suma acestei serii de puteri în intervalul $ (-1, 1) $. Derivăm
  termen cu termen și obținem:
  \[
    f'(x) = \sum (-1)^n x^{3n} = \frac{1}{1 + x^3},
  \]
  pentru $ |x| < 1 $, ca suma unei serii geometrice alternate.

  Rezultă:
  \[
    f(x) = \int \frac{dx}{1 + x^3} = \frac{1}{6} \ln \frac{(x + 1)^2}{x^2 - x + 1} %
    + \frac{1}{\sqrt{3}} \arctan{2x - 1}{\sqrt{3}} + c,
  \]
  pentru $ |x| < 1 $. Calculînd $ f(0) $, găsim $ c = \dfrac{\pi}{6 \sqrt{3}} $.
\item Se folosește seria pentru $ e^x $, din care obținem seria pentru
  $ (x + x^2)e^x $, pe care o derivăm termen cu termen.

  Pentru $ x = 1 $, se obține seria cerută, cu suma $ 5e $.
\item Descompunem seria în două, apoi folosim seria de puteri $ \sum n^2 x^n $,
  pe care o derivăm termen cu termen, pentru a obține seria pentru $ nx^{n-1} $,
  apoi seria pentru $ nx^n $.
\end{enumerate}

12. Studiați convergența seriilor de funcții:
\begin{enumerate}[(a)]
\item $ \sum \dfrac{n^2}{\sqrt{n!}} (x^n + x^{-n}) $, cu $ x \in \left[ \dfrac{1}{2}, 2 \right] $;
\item $ \sum \dfrac{nx}{1 + n^5 x^2}, x \in \RR $;
\item $ \sum \dfrac{\cos(3^n x)}{2^n}, x \in \RR $;
\item $ \sum x^n(1 - x), x \in [0, 1] $;
\item $ \sum \dfrac{(x + n)^2}{n^4}, x \in [0, 2] $;
\item $ \sum \dfrac{\ln(1 + nx)}{nx^n}, x > 0 $;
\item $ \sum \dfrac{xe^{-nx}}{\sqrt{n}}, x \geq 0 $;
\item $ \sum ne^{-nx}, x \geq 1 $.
\end{enumerate}

\emph{Indicații:} În fiecare caz, se găsește valoarea maximă a funcției $ | f_n | $,
pe care o notăm cu $ a_n $, apoi studiem convergența seriei numerice $ \sum a_n $.
%%% Local Variables:
%%% mode: latex
%%% TeX-master: "../m1-20-21"
%%% End:
